\appendix
\chapter{Annexe A - Business Model Canvas et Étude Économique}
\label{app:business}

Ce mémoire s'inscrivant dans le cadre de l'\textbf{Arrêté interministériel n\textsuperscript{o}~1275} (mécanisme «~Un diplôme, une Startup~»), cette annexe présente l'étude de faisabilité économique et le plan de développement commercial pour le projet de startup \textit{RohWinBghit}.

\section{Description et marché cible}

\subsection{Introduction}

La mobilité interurbaine en Algérie constitue un défi majeur. Les citoyens effectuant des trajets inter-wilayas sont souvent confrontés à des moyens de transport traditionnels (taxis collectifs, bus) qui souffrent d'horaires rigides, d'une asymétrie de l'information et d'une qualité de service variable. Parallèlement, de nombreux conducteurs voyagent seuls et supportent l'intégralité des coûts de leurs trajets. \textit{RohWinBghit} résout cette problématique en numérisant la mise en relation entre passagers et conducteurs, instaurant ainsi un cadre de confiance numérique permettant de fluidifier et de sécuriser la mobilité inter-wilayas.

\subsection{Description}

\textit{RohWinBghit} est une plateforme logicielle complète dédiée au covoiturage inter-wilayas. Ses fonctionnalités principales s'articulent autour de deux composantes :
\begin{itemize}
    \item \textbf{L'application mobile multiplateforme :} Destinée aux passagers et aux conducteurs, elle permet la recherche intuitive de trajets, la publication d'offres de covoiturage, la réservation de places avec billet QR, le suivi GPS en temps réel et la messagerie instantanée. 
    \item \textbf{Le panneau d'administration web :} Une interface de supervision dédiée aux administrateurs pour la modération des utilisateurs, la revue manuelle des dossiers KYC litigieux et le suivi des transactions financières.
\end{itemize}

\subsection{Marché cible}

Le marché potentiel en Algérie est particulièrement vaste. L'Algérie dispose d'une importante population étudiante répartie sur l'ensemble du territoire national, dont une part significative étudie en dehors de sa wilaya d'origine et nécessite des déplacements fréquents (pendulaires hebdomadaires). De plus, la forte adoption des smartphones et des services numériques favorise l'émergence de nouvelles solutions de mobilité, la population --- majoritairement jeune et connectée --- étant structurellement prête à adopter des solutions de mobilité dématérialisées.

\subsection{Analyse de la demande}

L'Algérie compte une forte demande pour des services de transport abordables, rapides et pratiques. Les alternatives existantes, telles que les taxis intercités, sont souvent perçues comme onéreuses, tandis que les groupes de covoiturage informels sur les réseaux sociaux manquent cruellement de mécanismes de sécurité et de traçabilité. Les usagers réclament une solution qui apporte des réponses concrètes à ces lacunes : des prix justes grâce à la mutualisation des frais, une vérification d'identité rigoureuse pour lever le frein de la peur, et un moyen de réservation confirmé à l'avance pour planifier les longs trajets en toute sérénité.

\section{Aspects d'innovation}

\subsection{Introduction}

L'innovation est essentielle pour assurer la compétitivité et la viabilité d'une entreprise technologique sur un marché émergent. Pour se différencier nettement de l'existant, \textit{RohWinBghit} intègre des innovations techniques et fonctionnelles avancées pour le marché algérien du transport inter-wilayas.

\subsection{Différence avec la concurrence}

En comparant notre produit aux solutions concurrentes, souvent limitées à de simples annuaires informels, notre plateforme présente des avantages technologiques décisifs :
\begin{itemize}
    \item \textbf{KYC biométrique automatisé :} Un pipeline d'intelligence artificielle combinant extraction OCR, comparaison faciale (ArcFace) et détection de vivacité vise à authentifier les profils, renforçant ainsi la confiance au sein de la plateforme.
    \item \textbf{Tarification dynamique :} Un algorithme de \textit{Surge Pricing} qui ajuste les tarifs en fonction de l'offre et de la demande spatio-temporelle, assurant un équilibre de marché équitable.
    \item \textbf{Paiement électronique local :} L'intégration native (actuellement en environnement sandbox) des cartes CIB et Edahabia pour dématérialiser les transactions.
\end{itemize}

Afin de renforcer la sécurité de notre écosystème, nous avons mis en place des protocoles stricts :
\begin{itemize}
    \item \textbf{Chiffrement des données :} Utilisation du chiffrement AES-256-GCM pour les données sensibles (plaques d'immatriculation) lors de leur stockage.
    \item \textbf{Authentification robuste :} Mise en œuvre d'une connexion sans mot de passe via OTP (One-Time Password) et gestion sécurisée des sessions par jetons JWT.
    \item \textbf{Sécurité réseau et hébergement souverain :} Les données biométriques sont traitées et hébergées localement afin d'être en parfaite conformité avec la Loi 25-11 relative à la protection des données personnelles.
\end{itemize}

\subsection{Impact sur la société}

Le déploiement de \textit{RohWinBghit} génère un impact positif mesurable à trois niveaux :
\begin{itemize}
    \item \textbf{Économique :} Réduction substantielle du coût des déplacements pour les usagers (notamment les étudiants) et contribution à la réduction des coûts de déplacement des propriétaires de véhicules grâce au partage des frais de trajet.
    \item \textbf{Social :} Renforcement de la sécurité des voyageurs, désenclavement de certaines régions moins desservies par les transports publics, et réduction de la fatigue liée à l'attente incertaine en station.
    \item \textbf{Environnemental :} L'optimisation du taux de remplissage des véhicules permet de limiter le nombre de voitures en circulation, contribuant directement à la réduction des émissions de CO\textsubscript{2} sur les axes autoroutiers.
\end{itemize}

\section{Modèle économique}

\subsection{Introduction}

Le modèle économique est le socle de la stratégie de croissance de \textit{RohWinBghit}. Il définit les mécanismes de création de valeur et les canaux de monétisation permettant d'assurer la viabilité financière de la plateforme. Notre objectif est de structurer un modèle capable de s'adapter rapidement à la croissance (scalabilité) tout en maintenant des coûts de fonctionnement maîtrisés.

\subsection{Description}

Notre modèle économique repose sur une approche de type \textbf{B2B2C / Peer-to-Peer}, où la plateforme agit comme un tiers de confiance facilitant les échanges entre les conducteurs (qui fournissent le service de mobilité) et les passagers (qui le consomment). La pérennité de l'entreprise s'appuie sur la génération de revenus issus de commissions transactionnelles et d'abonnements à forte valeur ajoutée.

\subsubsection{Value proposition (Propositions de valeur)}

La proposition de valeur de notre produit se décline en plusieurs points fondamentaux :
\begin{itemize}
    \item \textbf{Pour les passagers :} 
    \begin{itemize}[label=$\circ$]
        \item Accès à des trajets inter-wilayas économiques et confortables.
        \item Sécurité renforcée par la vérification biométrique des conducteurs.
        \item Réservation rapide et numérisée avec billet QR et suivi GPS.
    \end{itemize}
    \item \textbf{Pour les conducteurs :}
    \begin{itemize}[label=$\circ$]
        \item Couverture facile des frais de carburant et d'entretien sur les longs trajets.
        \item Flexibilité totale pour choisir les horaires et les itinéraires.
        \item Outils avancés de gestion des réservations et des revenus via le tableau de bord.
    \end{itemize}
\end{itemize}

\subsubsection{Customer segments (Clients)}

Les segments de clients pour cette plateforme peuvent être divisés en deux grandes catégories interdépendantes :

\paragraph{Segments de clients passagers :}
\begin{itemize}
    \item \textbf{Les étudiants universitaires :} Effectuant régulièrement des trajets entre leur domicile et les pôles universitaires en fin de semaine.
    \item \textbf{Les travailleurs pendulaires :} Les professionnels dont le lieu de travail est éloigné de leur wilaya de résidence et nécessitant un transport fiable et récurrent.
    \item \textbf{Les voyageurs occasionnels :} Les personnes cherchant une alternative rapide pour des visites familiales ou des rendez-vous ponctuels dans une autre wilaya.
\end{itemize}

\paragraph{Segments de clients conducteurs :}
\begin{itemize}
    \item \textbf{Les automobilistes réguliers :} Les propriétaires de véhicules effectuant des trajets inter-wilayas hebdomadaires ou quotidiens et cherchant à réduire leurs frais de déplacement.
    \item \textbf{Les professionnels itinérants :} Les commerciaux ou indépendants qui parcourent de longues distances dans le cadre de leur métier.
\end{itemize}

\subsubsection{Customer relationships (Relations)}

Nous abordons la relation client par le prisme de l'automatisation et de la confiance numérique :
\begin{itemize}
    \item \textbf{Système automatisé (\textit{Self-service}) :} L'application offre un parcours autonome complet, de l'inscription à la réservation.
    \item \textbf{Transparence et confiance :} Le système de notation croisée entre passagers et conducteurs régule naturellement la qualité du service.
    \item \textbf{Support et assistance :} Une équipe de modérateurs réactive disponible pour gérer les litiges et accompagner les utilisateurs.
\end{itemize}

\subsubsection{Channels (Canaux)}

Notre stratégie de distribution et d'acquisition s'appuie sur une approche multi-canale :
\begin{itemize}
    \item \textbf{Canaux digitaux directs :} Google Play Store et Apple App Store.
    \item \textbf{Acquisition sociale :} Campagnes ciblées sur les réseaux sociaux (Facebook, Instagram, TikTok) et partenariats avec des créateurs de contenu algériens.
    \item \textbf{Canaux physiques (BTL) :} Campagnes de tractage et d'affichage au sein des campus universitaires et des cités U.
\end{itemize}

\subsubsection{Key partners (Partenaires Clés)}

La croissance de l'écosystème repose sur des alliances stratégiques :
\begin{itemize}
    \item \textbf{Partenaires technologiques :} Fournisseurs d'infrastructure cloud et de services numériques, services de cartographie (Mapbox) et opérateurs SMS (OTP).
    \item \textbf{Partenaires financiers :} La passerelle de paiement en ligne \textbf{Chargily Pay V2} pour le traitement unifié des transactions par cartes CIB et Edahabia.
    \item \textbf{Partenaires institutionnels :} Universités, incubateurs et organismes d'accompagnement à l'innovation pour le soutien au développement et à la valorisation du projet.
\end{itemize}

\subsubsection{Key activities (Activités clés)}

Les activités centrales requises pour opérer la plateforme incluent :
\begin{itemize}
    \item Le développement logiciel continu et la maintenance de l'infrastructure technologique.
    \item La supervision manuelle des dossiers KYC et l'audit continu de sécurité.
    \item Le marketing, la communication et l'acquisition de nouveaux utilisateurs.
    \item La gestion du service client et la résolution des litiges.
\end{itemize}

\subsubsection{Key resources (Ressources clés)}

Pour soutenir notre activité, nous nous appuyons sur les ressources suivantes :
\begin{itemize}
    \item \textbf{Ressources technologiques :} La plateforme logicielle, l'infrastructure cloud et le pipeline de vérification biométrique intégrant plusieurs modèles d'intelligence artificielle.
    \item \textbf{Ressources humaines :} Une équipe fondatrice aux compétences pluridisciplinaires (développeurs, marketing, service client, direction).
    \item \textbf{Ressources immatérielles :} La marque déposée (INAPI), le code source scellé (ONDA) et la base croissante d'utilisateurs.
\end{itemize}

\subsubsection{Cost structure (Coûts)}

La structuration financière de \textit{RohWinBghit} se décompose entre les investissements initiaux (CAPEX) requis pour la mise en place juridique, technique et logistique de la startup, et les dépenses d'exploitation mensuelles régulières (OPEX) nécessaires au fonctionnement opérationnel de la plateforme.

\paragraph{Coûts initiaux (CAPEX) :}
Les investissements de démarrage englobent les dépenses indispensables avant le lancement commercial pour structurer l'entité légale, acquérir les équipements de l'équipe et obtenir les licences requises :
\begin{itemize}[leftmargin=1.5em]
    \item \textbf{Licences logicielles :} Acquisition du système d'exploitation Windows 11 pour les postes de l'équipe et abonnement à la suite collaborative Microsoft Office 365 Premium.
    \item \textbf{Coûts juridiques et administratifs :}
    \begin{itemize}
        \item \textbf{Enregistrement de l'entreprise :} Frais légaux d'immatriculation au Registre du Commerce (RC) et frais de notaire pour la rédaction des statuts.
        \item \textbf{Dépôt de marque et propriété intellectuelle :} Frais de dépôt de marque auprès de l'INAPI et de scellement du code source auprès de l'ONDA.
        \item \textbf{Conseils juridiques :} Honoraires d'avocats spécialisés pour la structuration légale de la startup et la rédaction des conditions générales d'utilisation (CGU) et des contrats de travail de l'équipe.
    \end{itemize}
    \item \textbf{Équipements et mobilier :}
    \begin{itemize}
        \item \textbf{Matériel informatique :} Acquisition de 7 ordinateurs de développement performants pour les membres de l'équipe.
        \item \textbf{Mobilier de bureau :} Acquisition de bureaux professionnels et de chaises ergonomiques (en distinguant le bureau du directeur de ceux du reste de l'équipe).
        \item \textbf{Équipement bureautique :} Achat d'imprimantes multifonctions pour la gestion administrative du bureau.
    \end{itemize}
\end{itemize}

Le tableau~\ref{tab:investissements} présente l'évaluation détaillée de ces investissements initiaux.

\begin{table}[H]
\centering
\small
\begin{tabularx}{\textwidth}{|X|r|c|r|}
\hline
\rowcolor{rohlight}
\textbf{Investissements initiaux} & \textbf{Montant (DZD)} & \textbf{Quantité} & \textbf{Total (DZD)} \\
\hline
Frais d'immatriculation au registre du commerce & 3\,700 & 1 & 3\,700 \\
\hline
Frais de dépôt de marque et de propriété intellectuelle (INAPI, ONDA) & 17\,000 & 1 & 17\,000 \\
\hline
Honoraires d'avocats (rédaction des statuts et des contrats) & 20\,000 & 1 & 20\,000 \\
\hline
Licences d'exploitation Windows 11 & 5\,000 & 7 & 35\,000 \\
\hline
Licence Microsoft Office 365 Premium & 25\,000 & 1 & 25\,000 \\
\hline
Ordinateurs de développement & 180\,000 & 7 & 1\,260\,000 \\
\hline
Bureaux professionnels (1 directeur + 6 collaborateurs) & 30\,000 + 45\,000 & 6 + 1 & 225\,000 \\
\hline
Chaises de bureau (1 directeur + 6 collaborateurs) & 18\,000 + 25\,000 & 6 + 1 & 133\,000 \\
\hline
Imprimantes multifonctions & 40\,000 & 2 & 80\,000 \\
\hline
\rowcolor{rohlight}
\textbf{Total des investissements initiaux (Somme 1)} & \textbf{--} & \textbf{--} & \textbf{1\,798\,700} \\
\hline
\end{tabularx}
  \caption{Évaluation des investissements initiaux nécessaires au lancement de RohWinBghit}
  \label{tab:investissements}
\end{table}

\paragraph{Dépenses régulières d'exploitation (OPEX) :}
Les charges récurrentes mensuelles assurent la continuité opérationnelle, la maintenance technique, la promotion marketing et la rémunération de l'équipe de 7 personnes :
\begin{itemize}[leftmargin=1.5em]
    \item \textbf{Salaires de l'équipe :} Rémunération mensuelle des 7 collaborateurs : 2 développeurs full-stack, 1 marketing manager, 2 agents d'acquisition campus et modération KYC, 1 team leader et 1 directeur général.
    \item \textbf{Loyer du local de bureau :} Location d'un appartement de type F3 implanté à Tlemcen, servant de bureau à l'équipe. Cette localisation régionale justifie le loyer modéré par rapport à la capitale.
    \item \textbf{Frais de marketing et publicité :}
    \begin{itemize}
        \item \textbf{Acquisition numérique :} Publicités payantes sur les moteurs de recherche et les médias sociaux.
        \item \textbf{Outils d'e-mailing :} Abonnement MailChimp pour les campagnes de fidélisation des conducteurs et passagers.
        \item \textbf{Partenariats d'influence :} Collaborations promotionnelles avec des influenceurs et blogueurs algériens pour la notoriété de la plateforme.
    \end{itemize}
    \item \textbf{Infrastructure cloud et services tiers :}
    \begin{itemize}
        \item \textbf{Hébergement cloud :} Serveurs VPS de production et base de données sécurisée.
        \item \textbf{Services GPU IA :} Calcul pour l'analyse de vivacité biométrique (vérification KYC des profils conducteurs).
        \item \textbf{API de géolocalisation :} Abonnement Mapbox pour le géo-routing, le calcul d'itinéraires et l'affichage des trajets inter-wilayas.
        \item \textbf{SMS OTP :} Envoi de messages d'authentification à usage unique pour la vérification des numéros de téléphone.
        \item \textbf{Passerelle de paiement :} Frais fixes d'abonnement à Chargily Pay pour l'intégration monétique locale.
        \item \textbf{Stockage médias (CDN) :} Cloudinary pour le stockage et la distribution des photos de profil, des documents KYC et des photos de véhicules.
    \end{itemize}
    \item \textbf{Logistique et divers :} Factures Internet, abonnements téléphoniques et frais de déplacements professionnels (visites campus, prospection).
\end{itemize}

Le tableau~\ref{tab:depenses} consolide l'ensemble de ces charges d'exploitation mensuelles.

\begin{table}[H]
\centering
\small
\begin{tabularx}{\textwidth}{|X|r|c|r|}
\hline
\rowcolor{rohlight}
\textbf{Désignation de la charge d'exploitation} & \textbf{Montant (DZD)} & \textbf{Quantité} & \textbf{Total mensuel (DZD)} \\
\hline
Salaire mensuel : Développeurs Full Stack & 65\,000 & 2 & 130\,000 \\
\hline
Salaire mensuel : Marketing manager & 80\,000 & 1 & 80\,000 \\
\hline
Salaire mensuel : Agents d'acquisition campus et modération KYC & 65\,000 & 2 & 130\,000 \\
\hline
Salaire mensuel : Team leader & 100\,000 & 1 & 100\,000 \\
\hline
Salaire mensuel : Directeur général & 130\,000 & 1 & 130\,000 \\
\hline
Loyer du local de bureau professionnel (F3) & 30\,000 & 1 & 30\,000 \\
\hline
Publicité sur les moteurs de recherche & 15\,000 & 1 & 15\,000 \\
\hline
Publicité sur les médias sociaux & 20\,000 & 1 & 20\,000 \\
\hline
Publicité et marketing d'e-mailing (MailChimp) & 16\,000 & 1 & 16\,000 \\
\hline
Partenariats avec des influenceurs locaux & 7\,500 & 3 & 22\,500 \\
\hline
Hébergement cloud et base de données (VPS Cloud) & 11\,000 & 1 & 11\,000 \\
\hline
Services GPU IA (analyse de vivacité biométrique KYC) & 10\,000 & 1 & 10\,000 \\
\hline
API de cartographie Mapbox (géo-routing des trajets) & 5\,000 & 1 & 5\,000 \\
\hline
Abonnements SMS OTP d'authentification & 2\,500 & 1 & 2\,500 \\
\hline
Passerelle de paiement Chargily Pay (frais fixes) & 2\,000 & 1 & 2\,000 \\
\hline
Cloudinary CDN (stockage photos profils, KYC, véhicules) & 4\,500 & 1 & 4\,500 \\
\hline
Abonnements Internet et téléphone & 4\,000 & 1 & 4\,000 \\
\hline
Frais de déplacements professionnels (visites campus) & 30\,000 & 1 & 30\,000 \\
\hline
\rowcolor{rohlight}
\textbf{Total des charges mensuelles (Somme 2)} & \textbf{--} & \textbf{--} & \textbf{742\,500} \\
\hline
\end{tabularx}
  \caption{Charges régulières mensuelles de fonctionnement de la plateforme}
  \label{tab:depenses}
\end{table}

\paragraph{Coût de lancement global :}
Le capital requis pour lancer la plateforme et financer le premier mois d'activité commerciale s'élève à :
\[ \text{Besoin global} = \text{CAPEX} + \text{OPEX}_{\text{M1}} = 1\,798\,700\text{ DZD} + 742\,500\text{ DZD} = 2\,541\,200\text{ DZD} \]

\subsubsection{Revenue Streams (Revenus)}

La plateforme \textit{RohWinBghit} adopte un modèle économique reposant sur plusieurs sources de revenus complémentaires, tout en maintenant un accès gratuit aux fonctionnalités essentielles pour les utilisateurs.

\paragraph{Commission sur les trajets (source principale) :}
Un prélèvement automatique de 12\,\% est effectué sur chaque transaction confirmée via la plateforme. Cette commission rémunère les services à valeur ajoutée proposés par \textit{RohWinBghit}, notamment la vérification d'identité (KYC biométrique), la sécurisation des échanges, la mise en relation des utilisateurs et l'accès aux fonctionnalités de réservation.

Sur la base d'un panier moyen estimé à 800~DZD par trajet, la plateforme perçoit une commission moyenne de 96~DZD par réservation. Dans un scénario de démarrage reposant sur environ 2\,813 trajets validés par mois, cette source génère un revenu estimé à 270\,000~DZD par mois, soit 3\,240\,000~DZD par an. Cette estimation repose sur un scénario prudent de démarrage et sera ajustée en fonction du rythme réel d'adoption de la plateforme.

Cette commission constitue la principale source de revenus de la plateforme dès son lancement.

\paragraph{Abonnement Premium Conducteur :}
\textit{RohWinBghit} prévoit également une offre Premium destinée aux conducteurs réguliers souhaitant bénéficier de fonctionnalités avancées telles qu'une meilleure visibilité dans les résultats de recherche, des statistiques détaillées sur leurs trajets et des avantages exclusifs proposés par la plateforme.

Le tarif de lancement est fixé à 500~DZD par mois. Cette source de revenus est appelée à se développer progressivement avec l'augmentation du nombre de conducteurs actifs.

\paragraph{Publicité locale ciblée (source complémentaire) :}
La plateforme prévoit à moyen terme l'intégration d'espaces publicitaires destinés aux commerces et services situés sur les principaux axes de déplacement (stations-service, restaurants, hôtels, garages automobiles et autres services de proximité).

Les revenus publicitaires sont considérés comme une source complémentaire et non comme un pilier du modèle économique initial. Dans un scénario prudent de démarrage, ils sont estimés à environ 30\,000~DZD par mois, soit 360\,000~DZD par an. Ces revenus pourront évoluer progressivement avec l'augmentation du nombre d'utilisateurs actifs et le développement de partenariats commerciaux.

\paragraph{Partenariats institutionnels et B2B :}
La plateforme prévoit le développement de partenariats avec des universités, des entreprises et des organismes de transport souhaitant faciliter les déplacements de leurs étudiants ou collaborateurs.

Ces collaborations pourront générer des revenus complémentaires à travers des conventions spécifiques et des offres adaptées aux besoins de chaque partenaire.

\subsubsection{Chiffre d'affaires projeté}

Dans un scénario prudent de démarrage, les revenus récurrents estimés sont les suivants :
\begin{itemize}[leftmargin=1.5em]
    \item Commission sur les trajets : 270\,000~DZD/mois.
    \item Abonnement Premium Conducteur : 50\,000~DZD/mois.
    \item Publicité locale ciblée : 30\,000~DZD/mois.
\end{itemize}

Le chiffre d'affaires récurrent estimé s'élève ainsi à environ 350\,000~DZD par mois, soit 4\,200\,000~DZD par an, hors revenus issus des partenariats B2B.

\subsubsection{Seuil de rentabilité opérationnelle}

Durant la phase de lancement, les revenus récurrents de démarrage (350\,000~DZD/mois) demeurent inférieurs aux charges d'exploitation mensuelles fixes (742\,500~DZD/mois), ce qui dégage un besoin de financement mensuel temporaire (\textit{burn rate}) de \textbf{392\,500~DZD}, couvert par les fonds propres des fondateurs, les dispositifs d'accompagnement à l'innovation ou les mécanismes de soutien aux startups.

Le seuil de rentabilité opérationnelle (point mort) sera progressivement atteint à moyen terme grâce à l'augmentation du volume de réservations (visant environ 6\,250 trajets mensuels pour couvrir le point mort), à la croissance du nombre de conducteurs Premium et au développement des partenariats institutionnels et commerciaux.

% ===== Branding Colors for BMC =====
\definecolor{bmcgreen}{RGB}{27,94,59}        % Deep Green #1B5E3B
\definecolor{bmclightgreen}{RGB}{232,245,238}% Soft Green #E8F5EE
\definecolor{bmclight}{RGB}{247,245,240}     % Light Beige #F7F5F0
\definecolor{bmcblue}{RGB}{27,94,59}         % Green synonym for existing headers
\definecolor{bmcgold}{RGB}{201,168,76}       % Gold #C9A84C

% Custom Box Styles for the Landscape Grid
\newtcolorbox{bmcblock}[2][]{
  enhanced,
  box align=top,
  colback=bmclight,
  colframe=bmcgreen,
  boxrule=1.2pt,
  arc=6pt,
  drop fuzzy shadow=black!12!white,
  title=#2,
  coltitle=white,
  colbacktitle=bmcgreen,
  fonttitle=\bfseries\sffamily\small,
  left=6pt, right=6pt, top=6pt, bottom=6pt,
  valign=top,
  #1
}

\newtcolorbox{bmchighlight}[2][]{
  enhanced,
  box align=top,
  colback=bmclightgreen,
  colframe=bmcgold,
  boxrule=1.8pt,
  arc=6pt,
  drop fuzzy shadow=bmcgreen!20!black,
  title=#2,
  coltitle=white,
  colbacktitle=bmcgold,
  fonttitle=\bfseries\sffamily\small,
  left=6pt, right=6pt, top=6pt, bottom=6pt,
  valign=top,
  #1
}


\clearpage % Force le placement de tous les flottants avant le changement de géométrie
\newgeometry{margin=1cm}%
\begin{landscape}
\thispagestyle{empty} % Masquer en-têtes et pieds de page pour une mise en page plein écran propre
\centering
\vspace*{\fill}
{\color{bmcgreen}\Large\bfseries\sffamily Business Model Canvas -- RohWinBghit}\\[0.4cm]
\renewcommand{\arraystretch}{1.0}%
\begin{tabularx}{\linewidth}{@{}XXXXX@{}}%
\begin{bmcblock}[height=10.6cm]{\faHandshake~Partenaires Clés}%
\scriptsize
\begin{itemize}[leftmargin=1.0em, itemsep=4pt, parsep=0pt]
  \item \textbf{Prototype et infrastructures} : Fournisseurs d'hébergement cloud, services de géolocalisation et cartographie, et passerelle monétique.
  \item \textbf{Soutien académique} : Universités et incubateurs technologiques nationaux.
  \item \textbf{Acteurs du transport public} : Partenariats stratégiques envisagés avec les opérateurs de transport.
  \item \textbf{Assurances} : Collaborations potentielles avec les compagnies d'assurance.
\end{itemize}
\end{bmcblock} &%
\begin{bmcblock}[height=4.5cm]{\faCogs~Activités Clés}%
\scriptsize
\begin{itemize}[leftmargin=1.0em, itemsep=2pt, parsep=0pt]
  \item \textbf{Développement continu} des applications mobiles et de la plateforme web.
  \item \textbf{Supervision KYC} et modération des documents d'identité.
  \item \textbf{Gestion des portefeuilles} et traitement des flux transactionnels.
\end{itemize}
\end{bmcblock}%
\vspace{0.3cm}%
\begin{bmcblock}[height=5.8cm]{\faLaptop~Ressources Clés}%
\scriptsize
\begin{itemize}[leftmargin=1.0em, itemsep=2pt, parsep=0pt]
  \item \textbf{Plateforme logicielle RohWinBghit} et base de données géographiques d'Algérie.
  \item \textbf{Modules de vision par ordinateur} formés pour la reconnaissance biométrique.
  \item \textbf{Équipe d'ingénieurs} et infrastructure cloud dédiée.
  \item \textbf{Propriété industrielle} (marques déposées).
\end{itemize}
\end{bmcblock} &%
\begin{bmchighlight}[height=10.6cm]{\faGem~Proposition de Valeur}%
\scriptsize
\begin{itemize}[leftmargin=1.0em, itemsep=4pt, parsep=0pt]
  \item \textbf{Covoiturage interurbain sécurisé} grâce à un contrôle d'identité automatique.
  \item \textbf{KYC biométrique} complet : Vérification documentaire et détection active de vivacité.
  \item \textbf{Paiement local flexible} : Support des cartes nationales CIB et Edahabia, et espèces.
  \item \textbf{Réservation numérique} traçable et billets sécurisés par codes QR.
  \item \textbf{Réduction des coûts} de trajet par partage des frais d'exploitation du véhicule.
\end{itemize}
\end{bmchighlight} &%
\begin{bmcblock}[height=4.5cm]{\faHeart~Relations Clients}%
\scriptsize
\begin{itemize}[leftmargin=1.0em, itemsep=2pt, parsep=0pt]
  \item \textbf{Confiance réciproque} : Profils certifiés et système d'évaluation mutuelle.
  \item \textbf{Support client} : Support et système de signalement des litiges intégré.
  \item \textbf{Accompagnement} des utilisateurs dans la phase KYC.
\end{itemize}
\end{bmcblock}%
\vspace{0.3cm}%
\begin{bmcblock}[height=5.8cm]{\faBullhorn~Canaux}%
\scriptsize
\begin{itemize}[leftmargin=1.0em, itemsep=2pt, parsep=0pt]
  \item \textbf{Magasins d'applications} mobiles d'iOS et d'Android.
  \item \textbf{Site internet} vitrine et console d'administration.
  \item \textbf{Réseaux sociaux} et campagnes de communication.
  \item \textbf{Campus universitaires} et programmes d'ambassadeurs.
\end{itemize}
\end{bmcblock} &%
\begin{bmcblock}[height=10.6cm]{\faUsers~Segments Clients}%
\scriptsize
\begin{itemize}[leftmargin=1.0em, itemsep=4pt, parsep=0pt]
  \item \textbf{Étudiants universitaires} effectuant des navettes inter-wilayas.
  \item \textbf{Travailleurs pendulaires} nécessitant des déplacements réguliers.
  \item \textbf{Conducteurs particuliers} cherchant à réduire leurs frais de déplacement.
  \item \textbf{Voyageurs occasionnels} se déplaçant pour raisons familiales ou loisirs.
\end{itemize}
\end{bmcblock} \\
\end{tabularx}

\vspace{0.4cm}

\begin{tabularx}{\linewidth}{@{}XX@{}}%
\begin{bmcblock}[height=4.2cm]{\faCoins~Structure de Coûts}%
\scriptsize
\begin{itemize}[leftmargin=1.5em, itemsep=1pt, parsep=0pt, topsep=2pt, partopsep=0pt]
  \item \textbf{Infrastructure \& DevOps} (\textasciitilde~28\,500~DZD/mois)
  \item \textbf{Services tiers} (\textasciitilde~9\,500~DZD/mois)
  \item \textbf{Ressources humaines} (\textasciitilde~570\,000~DZD/mois)
  \item \textbf{Marketing \& Acquisition} (\textasciitilde~73\,500~DZD/mois)
\end{itemize}
\vspace*{\fill}
\hrule
\vspace{0.1cm}
\centering
\begin{tabular}{l@{~:~}l@{\qquad\qquad}l@{~:~}l}
\textcolor{bmcgreen}{\textbf{CAPEX}} & \textbf{1\,798\,700~DZD} &
\textcolor{bmcgreen}{\textbf{OPEX}} & \textbf{742\,500~DZD / mois}
\end{tabular}
\end{bmcblock} &%
\begin{bmcblock}[height=4.2cm]{\faMoneyBillWave~Sources de Revenus}%
\scriptsize
\begin{itemize}[leftmargin=1.5em, itemsep=1pt, parsep=0pt, topsep=2pt, partopsep=0pt]
  \item \textbf{Commission sur les trajets} (\textasciitilde~270\,000~DZD/mois)
  \item \textbf{Abonnement Premium Conducteur} (\textasciitilde~50\,000~DZD/mois)
  \item \textbf{Publicité locale ciblée} (\textasciitilde~30\,000~DZD/mois)
  \item \textbf{Partenariats B2B} (revenus variables)
\end{itemize}
\vspace*{\fill}
\hrule
\vspace{0.1cm}
\centering
\begin{tabular}{l@{~:~}l@{\qquad\qquad}l@{~:~}l}
\textcolor{bmcgreen}{\textbf{CA estimé}} & \textbf{350\,000~DZD / mois} &
\textcolor{bmcgreen}{\textbf{CA estimé}} & \textbf{4\,200\,000~DZD / an}
\end{tabular}
\end{bmcblock} \\
\end{tabularx}

\vspace*{\fill}
\end{landscape}
\restoregeometry

Le Business Model Canvas présente une vue synthétique du modèle économique afin de faciliter la compréhension stratégique du projet. Les estimations financières détaillées, les hypothèses de coûts, les projections de revenus et l'analyse de rentabilité sont présentées dans la section Étude technico-économique et faisabilité du projet. Cette séparation permet de préserver la lisibilité du BMC tout en fournissant une analyse financière rigoureuse et cohérente avec les objectifs du dispositif « Un diplôme, une Startup » (Arrêté interministériel n\textsuperscript{o} 1275).
